\documentclass[a4paper]{article}
\usepackage{amsmath}
\usepackage{amssymb}
\usepackage{amsfonts}
\usepackage{bm}
\usepackage{geometry}
\geometry{left=25mm,right=25mm,top=25mm,bottom=25mm}
\usepackage{hyperref}

\title{Double Spacetime in Projective Geometric Algebra:\\
Duality, Null Translations, and Motors}

\author{https://github.com/hypernumbernet}
\date{\today}

\begin{document}

\maketitle

\begin{abstract}
Double spacetime theory treats gravity as torsional mismatch between two Minkowski frames carried by each massive particle. Those frames live in the biquaternion algebra \(\mathrm{Cl}(3,1)\). What they do not yet contain is translation. Projective geometric algebra \(G(3,1,1)\) supplies the missing piece by adjoining a single null vector \(e_{4}\). Translations become bivectors, on the same footing as boosts and rotations. The ten bivectors are a basis of the bivector space, which splits as the direct sum of three boosts, three rotations, and four translations. Under the commutator that space is the Poincar\'{e} algebra \(\mathfrak{so}(3,1)\ltimes\mathbb{R}^{3,1}\).

Three statements organise the construction. Duality \(X\mapsto Xi\) is a complex structure on the Lorentz sector: every boost is the dual of a rotation, and the sign that separates \(\mathfrak{so}(3,1)\) from \(\mathfrak{so}(4)\) is \(i^{2}=-1\). The geometric object that unites torsion with translation is the ordered product \(M=RT\) of a torsional rotor and a null translator, not a single Banach exponential. On a coordinate chart the exterior Schwarzschild geometry is a radial boost of that motor. The finite-angle scalar \(J\) that measures mismatch in parameter space is not the teleparallel field density \(T\).
\end{abstract}

\section{Introduction}
\label{sec:introduction}

The double spacetime theory refuses a shared continuum. Each massive particle carries a pair of Minkowski frames---the usual frame, in which we measure, and a dual frame, time-reversed relative to the first, with the same interval. Gravity is the torsion of that pair~\cite{Hestenes2015,DoranLasenby2003}. The algebra that holds the pair is the sixteen-dimensional biquaternion algebra, canonically isomorphic to \(\mathrm{Cl}(3,1)\). Boosts of the usual sector and rotations of the dual sector are generated by six bivectors. Their mismatch is a single algebraic object, and that object is enough to recover the exterior of a spherical mass without taking curvature as a primitive.

Six bivectors are not ten. The Lorentz algebra \(\mathfrak{so}(3,1)\) accounts for orientation; it does not account for displacement. To move a particle from one event to another, the algebra must also generate translations. The natural home for that generator, in geometric algebra, is a null vector: a direction that squares to zero and therefore produces, when wedged with the four Minkowski generators, four bivectors that themselves square to zero. This is the standard convention of projective geometric algebra \(G(3,1,1)\)~\cite{Gunn2017,Lengyel2024,DeKeninckDorst2019,GunnDeKeninck2020}, going back to Selig's treatment of rigid motions in degenerate Clifford algebras~\cite{Selig2005}. Adjoining one null vector \(e_{4}\) to \(\mathrm{Cl}(3,1)\) lifts the double spacetime into a thirty-two-dimensional algebra whose bivector space is ten-dimensional. Those ten bivectors are three boosts, three rotations, and four translations.

The lift is not a bookkeeping device. Under the commutator the ten-dimensional span closes, the translations form an abelian ideal, and the Lorentz sector acts on that ideal by its vector representation. The result is the Poincar\'{e} algebra \(\mathfrak{so}(3,1)\ltimes\mathbb{R}^{3,1}\) realised inside \(G(3,1,1)\)~\cite{Weinberg1995,Tung1985}. Duality, already present in the parent theory as right multiplication by the pseudoscalar \(i\), becomes a complex structure on the Lorentz factor: boosts are dualised rotations, and \(i^{2}=-1\) forces the boost--boost sign that distinguishes Lorentz from Euclidean. Translations, by contrast, are sent out of grade two. Duality organises orientation; it does not organise displacement.

Once both sectors are in hand, a motor is an ordered product \(M=RT\). The torsional factor \(R\) is the exponential of the six-dimensional mismatch; the translational factor \(T\) is linear, because a null bivector squares to zero. Unitarity \(M\widetilde{M}=1\) follows from the reversal-odd character of every bivector, which is why translations must be bivectors rather than vectors. When torsion and translation fail to commute, \(M\) is not the exponential of their sum. The geometric object is the product.

The same motor, restricted to a radial boost, recovers the diagonal Schwarzschild coframe. What it does not do is identify the parameter-space scalar \(J\) with the teleparallel torsion density \(T\). One is a finite angle; the other is a derivative of the coframe. Confusing the two collapses a constant boost into a curved geometry, or a pure translation into a non-vanishing field.

The paper follows that order. Section~\ref{sec:adjunction} adjoins \(e_{4}\) and classifies the ten generators. Section~\ref{sec:duality} shows that duality determines the Lorentz algebra. Section~\ref{sec:motors} constructs the motor and the quadratic diagnostic \(J^{(5)}\). Section~\ref{sec:gravity} reads the Schwarzschild exterior as a radial boost. Section~\ref{sec:conclusions} returns to what the algebra has settled and what it has not. The algebraic identities of those sections are machine-checked in the Lean~4 companion \texttt{DstDiophantine}; Appendix~\ref{app:lean} records the corresponding statements.

\section{Adjoining a null direction}
\label{sec:adjunction}

The parent theory identifies the four vector generators of \(\mathrm{Cl}(3,1)\) with biquaternion elements
\[
e_{0}=j,\qquad e_{1}=kI,\qquad e_{2}=kJ,\qquad e_{3}=kK,
\]
of Minkowski signature \((-,+,+,+)\):
\[
e_{0}^{2}=-1,\qquad e_{1}^{2}=e_{2}^{2}=e_{3}^{2}=+1,\qquad \{e_{\mu},e_{\nu}\}=0\quad(\mu\neq\nu).
\]
Projective geometric algebra \(G(3,1,1)\) is obtained by adjoining one further generator \(e_{4}\) with
\begin{equation}
e_{4}^{2}=0,\qquad \{e_{4},e_{\mu}\}=0\quad(\mu=0,1,2,3).
\label{eq:e4-defs}
\end{equation}
The full Clifford algebra is \(2^{5}=32\)-dimensional. The original sixteen-dimensional biquaternion algebra is the subalgebra generated by \(\{e_{0},e_{1},e_{2},e_{3}\}\).

The dual map of the parent theory is right multiplication by the biquaternionic pseudoscalar~\cite{ConwaySmith2003}
\[
i \;=\; e_{0}e_{1}e_{2}e_{3} \;=\; j\,(kI)(kJ)(kK).
\]
Because \(e_{4}\) anticommutes with each of the four Minkowski vectors, it commutes with their product:
\[
e_{4}\,i \;=\; e_{4}\,e_{0}e_{1}e_{2}e_{3} \;=\; (-1)^{4}\,e_{0}e_{1}e_{2}e_{3}\,e_{4} \;=\; i\,e_{4}.
\]
Duality therefore continues to act on \(\mathrm{Cl}(3,1)\) exactly as before, while the new null direction is left intact. The lift does not disturb the time-reversing character of \(X\mapsto Xi\).

The bivector space of \(G(3,1,1)\) has dimension \(\binom{5}{2}=10\). Those ten elements fall into three geometric classes, distinguished by the square.

The three \emph{hyperbolic} generators square to \(+1\) and span the boosts of \(\mathfrak{so}(3,1)\):
\[
iI \;=\; e_{0}e_{1},\qquad iJ \;=\; e_{0}e_{2},\qquad iK \;=\; e_{0}e_{3}.
\]
The three \emph{cyclic} generators square to \(-1\) and span the rotations of the same Lorentz copy:
\[
I \;=\; e_{3}e_{2},\qquad J \;=\; e_{1}e_{3},\qquad K \;=\; e_{2}e_{1}.
\]
The four \emph{null} generators are the wedges of \(e_{4}\) with the Minkowski basis,
\begin{align*}
N_{0} &= e_{4}\wedge e_{0} = e_{4}\,j, &
N_{1} &= e_{4}\wedge e_{1} = e_{4}\,(kI),\\
N_{2} &= e_{4}\wedge e_{2} = e_{4}\,(kJ), &
N_{3} &= e_{4}\wedge e_{3} = e_{4}\,(kK).
\end{align*}
From~\eqref{eq:e4-defs} one has, for every pair \(\mu,\nu\in\{0,1,2,3\}\),
\begin{equation}
N_{\mu}N_{\nu} \;=\; (e_{4}e_{\mu})(e_{4}e_{\nu}) \;=\; -\,e_{4}^{2}\,e_{\mu}e_{\nu} \;=\; 0.
\label{eq:null-strong-vanishing}
\end{equation}
Each \(N_{\mu}\) squares to zero, and every mixed product vanishes as well. The four null bivectors commute; they will be the translations.

A light-like translation is not a constraint on the generators. A general Minkowski vector \(L=ct\,e_{0}+x\,e_{1}+y\,e_{2}+z\,e_{3}\) is null when \(L^{2}=0\). The translation along \(L\) is the bivector \(e_{4}\wedge L\), a linear combination of the four \(N_{\mu}\), and it squares to zero by~\eqref{eq:null-strong-vanishing} whether or not \(L\) itself is null. Light-likeness lives on the four parameters \(\lambda^{\mu}=(ct,x,y,z)\), as the cone \(\eta_{\mu\nu}\lambda^{\mu}\lambda^{\nu}=0\). The translation sector remains four-dimensional.

\begin{table}[htbp]
\centering
\caption{The ten bivector generators of \(G(3,1,1)\).}
\label{tab:generators}
\begin{tabular}{llcl}
\hline
Class      & Generators & Count & Square \\
\hline
Hyperbolic & \(B_{1}^{+}=iI,\; B_{2}^{+}=iJ,\; B_{3}^{+}=iK\) & 3 & \(+1\) \\
Cyclic     & \(B_{1}^{-}=I,\; B_{2}^{-}=J,\; B_{3}^{-}=K\)     & 3 & \(-1\) \\
Null       & \(N_{0}=e_{4}j,\; N_{1}=e_{4}(kI),\; N_{2}=e_{4}(kJ),\; N_{3}=e_{4}(kK)\) & 4 & \(0\) \\
\hline
\end{tabular}
\end{table}

These ten elements are a basis of the bivector space. The three classes therefore span complementary subspaces, of dimensions three, three, and four, and every bivector is uniquely a boost plus a rotation plus a translation. The binomial count \(\binom{5}{2}=10\) is the dimension of this direct sum.

\section{Duality determines the Lorentz algebra}
\label{sec:duality}

The six hyperbolic and cyclic generators close among themselves. Same-axis boosts and rotations commute. Adjacent rotations reproduce the remaining rotation; adjacent boost--rotation pairs reproduce the remaining boost; adjacent boosts reproduce the remaining rotation with a reversed sign:
\begin{align*}
[B^{-}_{a},B^{-}_{a+1}] &= 2B^{-}_{a+2},\\
[B^{+}_{a},B^{-}_{a+1}] &= 2B^{+}_{a+2},\\
[B^{+}_{a},B^{+}_{a+1}] &= -2B^{-}_{a+2}.
\end{align*}
Indices run cyclically over the three axes. These are the structure constants of \(\mathfrak{so}(3,1)\) in the rotation--boost presentation \([J,J]\sim J\), \([J,K]\sim K\), \([K,K]\sim -J\). The minus sign in the last line is what separates the Lorentz algebra from compact \(\mathfrak{so}(4)\), where boosts would bracket to \(+J\).

The four null generators are moved among themselves by that Lorentz action. A radial boost exchanges time-translation with the spatial translation along its own axis and fixes the other two; a rotation turns the two translations in its plane into each other and fixes the rest. Representative brackets are
\begin{align*}
[B^{+}_{0},N_{0}]&=2N_{1},&
[B^{+}_{0},N_{1}]&=2N_{0},&
[B^{+}_{0},N_{2}]&=0,\\
[B^{-}_{0},N_{2}]&=2N_{3},&
[B^{-}_{0},N_{3}]&=-2N_{2},&
[B^{-}_{0},N_{0}]&=0.
\end{align*}
Every such bracket remains null. Null generators commute with one another, by~\eqref{eq:null-strong-vanishing}.

Write \(\mathfrak{l}\) for the six-dimensional span of boosts and rotations, and \(\mathfrak{n}\) for the four-dimensional null span. Their sum is the direct sum just obtained, and it is the whole bivector space. Bilinearity of the commutator lifts the generator identities to
\[
[\mathfrak{l},\mathfrak{l}]\subseteq\mathfrak{l},\qquad
[\mathfrak{l},\mathfrak{n}]\subseteq\mathfrak{n},\qquad
[\mathfrak{n},\mathfrak{n}]=0.
\]
Thus \(\mathfrak{l}\) and \(\mathfrak{p}:=\mathfrak{l}\oplus\mathfrak{n}\) are Lie subalgebras of \(G(3,1,1)\), and \(\mathfrak{n}\) is an abelian ideal of \(\mathfrak{p}\). The quotient \(\mathfrak{p}/\mathfrak{n}\) is \(\mathfrak{l}\), and \(\mathfrak{l}\) acts on \(\mathfrak{n}\) by the vector representation just recorded. This is the semidirect product \(\mathfrak{so}(3,1)\ltimes\mathbb{R}^{3,1}\) as an algebraic fact about the ten bivectors. The action is faithful in the sense that motors require: \([B^{+}_{0},N_{0}]=2N_{1}\neq 0\), so torsion and translation do not commute in general.

The internal structure of \(\mathfrak{l}\) is duality. The map \(X\mapsto Xi\) squares to \(-1\), because \(i^{2}=-1\). The pseudoscalar \(i\) anticommutes with each vector \(e_{\mu}\) and therefore commutes with every product of two of them: \(i\) is central on \(\mathfrak{l}\). Together these say that right multiplication by \(i\) is a complex structure on the Lorentz sector, and that \(\mathfrak{l}\) is a module over \(\mathbb{R}[i]\cong\mathbb{C}\).

On the generators the complex structure is exact:
\[
B^{-}_{a}\,i \;=\; B^{+}_{a},\qquad
B^{+}_{a}\,i \;=\; -B^{-}_{a}\qquad(a=0,1,2).
\]
Every boost is the dual of the corresponding rotation. If \(\mathfrak{r}\) is the three-dimensional span of the rotation generators, the boost sector is \(\mathfrak{r}\,i\), and
\[
\mathfrak{l} \;=\; \mathfrak{r}\;\oplus\;\mathfrak{r}\,i.
\]
The Lorentz algebra is the complexification of the rotation algebra. This is the classical isomorphism \(\mathfrak{so}(3,1)\cong\mathfrak{so}(3)\otimes\mathbb{C}\cong\mathfrak{sl}(2,\mathbb{C})\), written inside \(G(3,1,1)\) as the statement that boosts are dualised rotations.

The commutator respects the structure. Since \(i\) commutes with every element of \(\mathfrak{l}\),
\[
[Xi,Y] \;=\; [X,Y]\,i \;=\; [X,Yi],\qquad
[Xi,Yi] \;=\; [X,Y]\,i^{2} \;=\; -[X,Y]
\qquad(X,Y\in\mathfrak{l}).
\]
The bracket is complex-bilinear. Consequently the entire Lorentz table is generated by the rotation table alone. From \([B^{-}_{a},B^{-}_{a+1}]=2B^{-}_{a+2}\) one obtains
\[
[B^{+}_{a},B^{-}_{a+1}] \;=\; [B^{-}_{a},B^{-}_{a+1}]\,i \;=\; 2B^{+}_{a+2},
\qquad
[B^{+}_{a},B^{+}_{a+1}] \;=\; -[B^{-}_{a},B^{-}_{a+1}] \;=\; -2B^{-}_{a+2}.
\]
The minus sign in the boost--boost bracket is \(i^{2}=-1\). In Euclidean signature the four-dimensional pseudoscalar squares to \(+1\), the same computation yields \([K,K]\sim +J\), and the complexification collapses to compact \(\mathfrak{so}(4)\cong\mathfrak{so}(3)\oplus\mathfrak{so}(3)\). The Lorentz signature of the double spacetime is encoded in the duality map itself.

The null sector behaves differently. Because \(e_{4}\) commutes with \(i\),
\[
N_{\mu}\,i \;=\; (e_{4}e_{\mu})\,i \;=\; e_{4}\,(e_{\mu}\,i),
\]
and \(e_{\mu}i\) is grade three in \(\mathrm{Cl}(3,1)\), so \(N_{\mu}i\) is grade four in \(G(3,1,1)\). Duality sends each translation out of the bivector grade. The null ideal is not dual-closed. The complex structure belongs to the Lorentz factor of \(\mathfrak{p}=\mathfrak{l}\oplus\mathfrak{n}\), not to the Poincar\'{e} span as a whole.

\section{Motors}
\label{sec:motors}

The extended mismatch is the sum of a torsional bivector and a translational bivector,
\begin{equation}
\Omega_{\mathrm{biv}}^{(5)}
= \underbrace{\sum_{a=1}^{3} \Bigl( \frac{\alpha_{a}}{2}\, B_{a}^{+} + \frac{\beta_{a}}{2}\, B_{a}^{-} \Bigr)}_{\Omega_{\mathrm{torsion}}}
\;+\;
\underbrace{\sum_{\mu=0}^{3} \frac{\lambda^{\mu}}{2}\, N_{\mu}}_{\Omega_{\mathrm{trans}}}.
\label{eq:Omega-decomp}
\end{equation}
The six coefficients \(\alpha_{a}\), \(\beta_{a}\) are the original mismatch of usual and dual sectors. The four coefficients \(\lambda^{\mu}\) are a Lorentz four-vector of translational mismatch. Light-like translations are the cone \(\eta_{\mu\nu}\lambda^{\mu}\lambda^{\nu}=0\).

Because every mixed product \(N_{\mu}N_{\nu}\) vanishes, the translational generator is nilpotent:
\[
\Omega_{\mathrm{trans}}^{2}
\;=\;
\frac{1}{4}\sum_{\mu,\nu}\lambda^{\mu}\lambda^{\nu}\,N_{\mu}N_{\nu}
\;=\; 0.
\]
Its exponential therefore truncates at first order,
\[
\exp(\Omega_{\mathrm{trans}}) \;=\; 1+\Omega_{\mathrm{trans}}.
\]
If \(T(p)=1+\Omega_{\mathrm{trans}}(p)\) and \(T(q)=1+\Omega_{\mathrm{trans}}(q)\), the product is \(T(p)T(q)=T(p+q)\). Translators form an abelian group under multiplication. A translation is a pure shift; it introduces neither curvature nor higher-order terms.

The geometric motor is the ordered product of the two exponentials,
\[
M \;:=\; R\,T,
\qquad
R \;:=\; \exp(\Omega_{\mathrm{torsion}}),\qquad
T \;:=\; 1+\Omega_{\mathrm{trans}}.
\]
The order matters. Cross-commutators \([B_{a}^{\pm},N_{\mu}]\) are generally nonzero, so \(\Omega_{\mathrm{torsion}}\) and \(\Omega_{\mathrm{trans}}\) need not commute, and \(M\) need not equal \(\exp(\Omega_{\mathrm{biv}}^{(5)})\). When they do commute, the two expressions coincide. When they do not---a radial boost against a time translation is the simplest example---the geometric object remains the product \(RT\). A pure radial boost conjugates translators to translators: the sandwich \(R(\varphi)\,T(p)\,\widetilde{R(\varphi)}\) is the translator whose coefficients are the Lorentz action of rapidity \(\varphi\) on the boost plane.

Unitarity of \(M\) is a consequence of grade. Reversal sends every bivector to its negative, \(\widetilde{B_{a}^{\pm}}=-B_{a}^{\pm}\) and \(\widetilde{N_{\mu}}=-N_{\mu}\). Had translations been taken as grade-one null vectors, they would have been reversal-even, and the argument below would fail. That is why projective geometric algebra generates translations by null \emph{bivectors}~\cite{Gunn2017,Lengyel2024,GunnDeKeninck2020}. For the rotor one has the standard cancellation
\[
R\,\widetilde{R}
\;=\;
\exp(\Omega_{\mathrm{torsion}})\,\exp(-\Omega_{\mathrm{torsion}})
\;=\; 1.
\]
For the translator, reversal-oddness together with nilpotence gives
\[
T\,\widetilde{T}
\;=\;
\bigl(1+\Omega_{\mathrm{trans}}\bigr)\bigl(1-\Omega_{\mathrm{trans}}\bigr)
\;=\;
1-\Omega_{\mathrm{trans}}^{2}
\;=\; 1.
\]
Reversal is an antihomomorphism, so
\[
M\,\widetilde{M}
\;=\;
R\,T\,\widetilde{T}\,\widetilde{R}
\;=\;
R\,\widetilde{R}
\;=\; 1.
\]
The motor is unitary whether or not the factors commute.

Duality acts on the torsional parameters as a quarter-turn. Using \(B^{-}_{a}i=B^{+}_{a}\) and \(B^{+}_{a}i=-B^{-}_{a}\),
\[
\Omega_{\mathrm{torsion}}(\alpha,\beta)\,i
\;=\;
\sum_{a}\Bigl(\frac{\beta_{a}}{2}\,B^{+}_{a}-\frac{\alpha_{a}}{2}\,B^{-}_{a}\Bigr)
\;=\;
\Omega_{\mathrm{torsion}}(\beta,-\alpha).
\]
Applied twice, the map reverses all six coefficients, in agreement with \(i^{2}=-1\). Writing \(z_{a}=\alpha_{a}+\mathrm{i}\beta_{a}\), duality is multiplication by \(-\mathrm{i}\). The torsional scalar
\[
J \;=\; \tfrac12\sum_{a}(\alpha_{a}^{2}-\beta_{a}^{2}) \;=\; \tfrac12\operatorname{Re}\sum_{a}z_{a}^{2}
\]
is the real part of the holomorphic quadratic; the unsigned mass \(\tfrac12\sum_{a}|z_{a}|^{2}\) is its Hermitian companion. Duality therefore reverses \(J\) and preserves the mass. The usual/dual swap of the parent theory, which exchanges \(\alpha_{a}\) with \(\beta_{a}\), differs from this quarter-turn only by the sign of the new rotation parameter; both reverse \(J\). Translational parameters must be tracked separately, because duality sends each \(N_{\mu}\) to grade four.

A quadratic diagnostic on the full ten-dimensional space cannot come from the Cartan--Killing form of \(\mathfrak{iso}(3,1)\) alone. That form vanishes on translations~\cite{FultonHarris1991,Humphreys1972}. Extending the Killing quadratic of the Lorentz sector by the Minkowski metric on the translation coefficients gives a non-degenerate pairing
\begin{equation}
J^{(5)}
\;:=\;
\tfrac{1}{16}\,B_{\mathrm{Killing}}(\Omega_{\mathrm{torsion}},\Omega_{\mathrm{torsion}})
\;+\;
\tfrac{1}{2}\,\eta_{\mu\nu}\lambda^{\mu}\lambda^{\nu}.
\label{eq:J5-invariant}
\end{equation}
Under the half-angle expansion of~\eqref{eq:Omega-decomp}, with \(B(B_{a}^{+},B_{a}^{+})=+8\) and \(B(B_{a}^{-},B_{a}^{-})=-8\), the first term equals \(\tfrac18\sum_{a}(\alpha_{a}^{2}-\beta_{a}^{2})\), which is one quarter of the discrete companion's \(J=\tfrac12\sum_{a}(\alpha_{a}^{2}-\beta_{a}^{2})\). The second term vanishes precisely on the light cone. The two contributions decouple because the Killing form of \(\mathfrak{iso}(3,1)\) vanishes between the semisimple and translation generators. The translational contribution is unbounded unless \(\lambda^{\mu}\) is constrained; the admissible ceiling \(|J|\le 3\pi^{2}/8\) of the discrete companion applies to the torsion part alone.

One geometric reading of the same algebra is worth recording, although it is not needed for the motor calculus. Planes in \(G(3,1,1)\) are three-blades. The antisymmetric part of their geometric product,
\[
A \mathbin{\underline{\vee}} B \;:=\; \tfrac12(AB-BA),
\]
takes values in grade two. On representative configurations, if \(P\) is a three-blade and \(D\) a dual-sector element aligned with the normal \(n\) of the plane of \(P\), the bracket is a null bivector along that normal,
\[
P \mathbin{\underline{\vee}} D \;=\; \kappa\,(e_{4}\wedge n).
\]
The dual spacetime then acts as a carrier of normals, and the bracket produces a translation along the normal. The appendix works one such example. A fully covariant calculus relating these brackets to \(\Omega_{\mathrm{biv}}^{(5)}\) is not developed here.

\section{Gravity on a chart}
\label{sec:gravity}

The parent theory aims to realise teleparallel gravity: geometry is carried by torsion rather than by a curvature tensor. In the projective extension the same idea is available on a fixed coordinate chart. Throughout we work with the Poincar\'{e} subalgebra of Section~\ref{sec:duality}. Variational equivalence of the teleparallel scalar with the Einstein--Hilbert action on a general manifold is taken from the TEGR literature~\cite{maluf2013} and is not re-derived.

A unitary motor acts on the reference coframe by sandwiching. Fix the \(\mathrm{Cl}(3,1)\) vector generators \(\{\mathring{e}^{a}\}_{a=0}^{3}\) inside \(G(3,1,1)\), and set
\begin{equation}
e^{a} \;:=\; \bigl\langle M\,\mathring{e}^{a}\,\widetilde{M}\bigr\rangle_{1}.
\label{eq:motor-tetrad}
\end{equation}
On a chart \((x^{\mu})\) the component fields \(e^{a}{}_{\mu}(x)\) induce the metric
\begin{equation}
g_{\mu\nu} \;=\; \eta_{ab}\, e^{a}{}_{\mu}\, e^{b}{}_{\nu},
\label{eq:induced-metric}
\end{equation}
with \(\eta=\mathrm{diag}(-1,+1,+1,+1)\). When \(M\) is a pure Lorentz rotor, the sandwich is a local Lorentz action on the reference vectors. Null translations enter the full motor \(M=RT\) but do not by themselves supply the classical Schwarzschild diagonal gauge; they remain available as translational mismatch measured by \(J^{(5)}\).

The Weitzenb\"{o}ck torsion of a coframe on the chart is the antisymmetrised derivative
\begin{equation}
T^{a}{}_{\mu\nu} \;=\; \partial_{\mu} e^{a}{}_{\nu} - \partial_{\nu} e^{a}{}_{\mu}.
\label{eq:weitzenbock-torsion}
\end{equation}
With inverse tetrad \(e_{a}{}^{\mu}\) and spacetime indices lowered by \(g_{\mu\nu}\), the teleparallel scalar is
\begin{equation}
T \;=\; \tfrac14 T_{\rho\mu\nu}T^{\rho\mu\nu} + \tfrac12 T_{\rho\mu\nu}T^{\nu\mu\rho} - T_{\rho}T^{\rho},
\label{eq:teleparallel-T}
\end{equation}
where \(T_{\rho}=T^{\sigma}{}_{\rho\sigma}\). No curvature tensor of a metric-compatible Levi-Civita connection is required: the teleparallel description stores geometry in \(T^{a}{}_{\mu\nu}\).

The exterior Schwarzschild geometry is a radial boost of this construction. Let \(r_{s}=2GM/c^{2}>0\) and work in the region \(r>r_{s}\). The diagonal coframe (units \(c=1\)) is
\begin{align}
e^{0} &= \sqrt{1-\frac{r_{s}}{r}}\,dt,\nonumber\\
e^{1} &= \bigl(1-\frac{r_{s}}{r}\bigr)^{-1/2}\,dr,\nonumber\\
e^{2} &= r\,d\theta,\nonumber\\
e^{3} &= r\sin\theta\,d\phi.
\label{eq:schwarzschild-tetrad}
\end{align}
Set \(A(r):=1-r_{s}/r\) and
\begin{equation}
\varphi(r) \;:=\; -\log\sqrt{A(r)} \;=\; -\tfrac12\log A(r)
\qquad(r>r_{s}).
\label{eq:radial-rapidity}
\end{equation}
Then \(e^{-\varphi}=\sqrt{A}\) and \(e^{\varphi}=A^{-1/2}\), so the \((t,r)\) diagonal coefficients are the boost scale factors of the pure hyperbolic rotor generated by \(B^{+}_{1}=e_{0}\wedge e_{1}\). The angular legs come from the spherical reference coframe. The corresponding torsional parameter is the pure boost \(\Omega_{\mathrm{torsion}}=(\varphi/2)B^{+}_{1}\), with vanishing null parameters in the diagonal gauge. The induced metric~\eqref{eq:induced-metric} is the Schwarzschild metric.

Three scalars now sit on the same chart, and they are not the same object. The finite-angle diagnostic
\[
J(\varphi) \;=\; \tfrac12\varphi(r)^{2}
\]
measures the size of the boost in parameter space. The field seed
\begin{equation}
J_{\mathrm{field}}(r) \;:=\; \tfrac12\bigl(\varphi'(r)\bigr)^{2},
\label{eq:J-field}
\end{equation}
with \(\varphi'=-r_{s}/(2r^{2}A)\), measures the radial rate of that boost. The teleparallel scalar~\eqref{eq:teleparallel-T} is a quadratic in the derivatives of the coframe. On this exterior chart it evaluates to
\[
T \;=\; \frac{2}{r^{2}}\frac{(1-\sqrt{A})^{2}}{\sqrt{A}},
\]
and with the TEGR radial current \(B^{r}=4(1-\sqrt{A})/r\) it is a pure radial divergence,
\[
T \;=\; -r^{-2}\,\partial_{r}(r^{2}B^{r}).
\]
The identity \(\partial_{r}\sqrt{A}=-\varphi'\sqrt{A}\) relates the boost torsion leg to \(\varphi'\), but it does not identify \(J\) or \(J_{\mathrm{field}}\) with \(T\).

A constant nonzero boost has \(J\neq 0\) while a flat chart has \(T=0\). A pure spatial translation on Minkowski space has \(J^{(5)}\neq 0\) and \(T=0\). Even the field seed fails to match \(\tfrac12 T\) at a generic exterior point: at \((r_{s},r)=(2,4)\), for instance, the two numbers disagree. Finite angle, rate of angle, and coframe density live on different rungs of the derivative ladder. A dictionary for a general smooth motor field \(M(x)=R(x)T(x)\) would have to be built from derivatives of \((\Omega_{\mathrm{torsion}},\lambda)\), not from finite angles, and would be expected to involve a total divergence. That dictionary is not constructed here.

Sandwiching with respect to the full degenerate metric of \(G(3,1,1)\) is likewise left aside. The identification of \(M\) with \(\exp(\Omega_{\mathrm{biv}}^{(5)})\) holds when the factors commute and fails in general otherwise.

\section{Conclusions}
\label{sec:conclusions}

Adjoining a null vector \(e_{4}\) to the biquaternion algebra of double spacetime realises translations as bivectors. The ten generators are a basis of the bivector space, and that space is the Poincar\'{e} algebra \(\mathfrak{so}(3,1)\ltimes\mathbb{R}^{3,1}\): the six torsional generators close, and the four null generators span an abelian ideal on which they act. Duality \(X\mapsto Xi\) is the complex structure of the Lorentz factor. Every boost is a dualised rotation, the bracket is complex-bilinear, and the sign that separates \(\mathfrak{so}(3,1)\) from \(\mathfrak{so}(4)\) is \(i^{2}=-1\). Translations are sent to grade four; duality does not close the null ideal.

The motor that unites the two sectors is the ordered product \(M=RT\). Nilpotence of the null generator truncates its exponential to a translator \(T=1+\Omega_{\mathrm{trans}}\). Reversal-oddness of every bivector gives \(M\widetilde{M}=1\). When torsion and translation fail to commute, the product is not a single exponential; it remains the geometric object. Duality turns the torsional parameters through a quarter-turn, reversing \(J\) and preserving the unsigned mass. The diagnostic \(J^{(5)}\) extends the Killing quadratic by the Minkowski pairing on translations; its translational term is unbounded off the light cone.

On a chart, sandwiching with \(M\) produces a tetrad. The exterior Schwarzschild coframe is a radial boost of that tetrad. The finite-angle scalar \(J\), the field seed \(J_{\mathrm{field}}\), and the teleparallel density \(T\) are three different readings of the same geometry; they are not pointwise identical. What remains open is a covariant bracket calculus for planes, a sandwich with the degenerate metric of \(G(3,1,1)\), and a field dictionary for a general motor. The variational equivalence of TEGR with Einstein--Hilbert gravity is the input of~\cite{maluf2013}.

\appendix

\section{A plane bracket producing a translation}
\label{app:bracket}

Planes in \(G(3,1,1)\) are three-blades. For two three-blades \(A\) and \(B\) the antisymmetric and symmetric parts of the geometric product are
\[
A \mathbin{\underline{\vee}} B \;=\; \tfrac12(AB-BA),\qquad
A \mathbin{\overline{\vee}} B \;=\; \tfrac12(AB+BA).
\]
These symbols are used only in this sense; they are not the anti-inner and anti-outer products of the PGA literature~\cite{Lengyel2024,DorstFontijneMann2007}. In a five-dimensional vector space the geometric product of two grade-three elements has grades \(0\), \(2\), and \(4\) only.

Take \(P=e_{4}\wedge e_{1}\wedge e_{2}\) and \(D=jI=-e_{0}e_{3}e_{2}\), a dual-sector trivector of the parent theory. Anticommuting generators through the product \(PD-DP\) reduces the bracket to a multiple of \(e_{4}\wedge e_{3}\). The direction \(e_{3}\) is the unit normal to the spatial plane \(e_{1}\wedge e_{2}\), so
\[
P \mathbin{\underline{\vee}} D \;=\; \kappa\, N_{3}
\]
for a real coefficient \(\kappa\). The right-hand side is a generator of translations along that normal. The dual element \(D\) has acted as the carrier of the normal direction.

A second, purely spatial example confirms the grade pattern. For \(A=e_{1}e_{2}e_{3}\) and \(B=e_{1}e_{2}e_{4}\) one finds \(AB=-e_{3}e_{4}\) and \(BA=+e_{3}e_{4}\), hence \(A\mathbin{\underline{\vee}}B=-e_{3}e_{4}\) and \(A\mathbin{\overline{\vee}}B=0\): a pure grade-two commutator, as expected of an antisymmetric product.

\section{Lean-checked identities}
\label{app:lean}

{\raggedright
\emergencystretch=3em

The construction of the paper is formalised in the Lean~4 project \texttt{DstDiophantine}. What follows is not a second proof, but a map from the three claims of the introduction to the theorems that check them. Proofs live in the named modules; they are not reproduced here. An abstract isomorphism with \(\mathfrak{iso}(3,1)\), a field dictionary for a general motor, and sandwiching against the full degenerate metric of \(G(3,1,1)\) are not claimed in the formalisation. Linear independence of the ten generators, and the direct sum of the three classes inside the bivector space, are \texttt{linearIndependent\_tenGen}, \texttt{poincareSpan\_eq\_bivectorGrade}, and \texttt{disjoint\_lorentz\_null} in \texttt{Algebra.BivectorBasis}.

\subsection{Null direction and the ten generators}

The algebra is the Clifford algebra of the quadratic form with weights \((-1,+1,+1,+1,0)\), in modules \texttt{Algebra.QuadraticForm} and \texttt{Algebra.PGA}. The identities \(e_{4}^{2}=0\) and \(\{e_{4},e_{\mu}\}=0\) are \texttt{PGA.e4\_sq} and \texttt{PGA.e4\_anticomm}. The ten bivectors of Table~\ref{tab:generators} are \texttt{hyperbolic}, \texttt{cyclic}, and \texttt{null} in \texttt{Algebra.Generators}. Squares and strong vanishing read
{\small
\begin{verbatim}
theorem hyperbolic_sq (a : Fin 3) : hyperbolic a * hyperbolic a = 1
theorem cyclic_sq (a : Fin 3) : cyclic a * cyclic a = -1
theorem null_mul_null (mu nu : Fin 4) : null mu * null nu = 0
theorem null_reverse (mu : Fin 4) : reverse (null mu) = -null mu
\end{verbatim}
}
The last line is why translations must be bivectors: reversal of a grade-one null vector would be even, and the unitarity argument of Section~\ref{sec:motors} would fail. The same module proves \texttt{hyperbolic\_reverse} and \texttt{cyclic\_reverse}.

\subsection{Duality and the Poincar\'{e} span}

\texttt{Algebra.LorentzLie} lifts the generator table to spans. Write \(\mathfrak{l}\) for \texttt{lorentzSpan} and \(\mathfrak{n}\) for \texttt{nullSpan}. The Lorentz span and the ten-dimensional span close under the commutator (\texttt{commutator\_mem\_lorentzSpan}, \texttt{commutator\_mem\_poincareSpan}). The null span is an abelian ideal (\texttt{commutator\_null\_null\_eq\_zero}, \texttt{commutator\_poincare\_null\_mem\_nullSpan}). These are the mathlib Lie subalgebras \texttt{lorentzLieSubalgebra} and \texttt{poincareLieSubalgebra}. The action is faithful in the sense motors require: \texttt{commutator\_hyperbolic0\_null0\_ne\_zero} is \([B^{+}_{0},N_{0}]\neq 0\).

Duality is right multiplication by the Cl(3,1) pseudoscalar. On the Lorentz span it is a complex structure, and the bracket is complex-bilinear:
{\small
\begin{verbatim}
theorem hyperbolic_eq_dual_cyclic (a : Fin 3) :
    hyperbolic a = dual (cyclic a)
theorem dual_dual (x : PGA) : dual (dual x) = -x
theorem commutator_dual_dual (x y : PGA)
    (hx : x Mem lorentzSpan) (hy : y Mem lorentzSpan) :
    commutator (dual x) (dual y) = -commutator x y
theorem commutator_hyperbolic_hyperbolic_cyclic (a b : Fin 3) :
    commutator (hyperbolic a) (hyperbolic b) =
      if a = b then 0
      else if b = a + 1 then -((2 : Real) smul cyclic (a + 2))
      else (2 : Real) smul cyclic (b + 2)
\end{verbatim}
}
The minus sign in the boost--boost table is \(i^{2}=-1\). Every Lorentz element is a rotation plus a dualised rotation (\texttt{mem\_lorentzSpan\_iff}). On torsion parameters duality is the quarter-turn \((\alpha,\beta)\mapsto(\beta,-\alpha)\) (\texttt{dual\_omegaTorsion}). It reverses \(J\) and preserves the unsigned mass (\texttt{J\_dualParams}, \texttt{mass\_dualParams}).

\subsection{Motors}

\texttt{Algebra.Motor} defines \(\Omega_{\mathrm{torsion}}\), \(\Omega_{\mathrm{trans}}\), and the ordered product \(RT\). Nilpotence of the translational generator truncates the Banach exponential, translators add, and the product is unitary:
{\small
\begin{verbatim}
theorem omegaTrans_sq (p : TransParams) :
    omegaTrans p * omegaTrans p = 0
theorem exp_omegaTrans (p : TransParams) :
    exp (omegaTrans p) = 1 + omegaTrans p
theorem expTrans_mul (p q : TransParams) :
    expTrans p * expTrans q =
      expTrans (fun mu => p.lambda mu + q.lambda mu)
theorem rotor_unitary (p : TorsionParams) :
    rotorTorsion p * reverse (rotorTorsion p) = 1
theorem motor_unitary (p : OmegaParams) :
    motor p * reverse (motor p) = 1
\end{verbatim}
}
When the factors commute, \(\exp(\Omega_{\mathrm{biv}}^{(5)})=RT\) (\texttt{exp\_omegaBiv\_eq\_motor\_of\_commute}). When they do not, the two differ: \texttt{exists\_omegaBiv\_ne\_motor} exhibits a radial boost against a time translation. Every torsion rotor, not only a pure radial boost, conjugates translators to translators (\texttt{Algebra.MotorGroup}, \texttt{exists\_sandwich\_rotorTorsion\_expTrans}). Closed forms \(\exp(tB^{+})=\cosh t+\sinh t\cdot B^{+}\) and \(\exp(tB^{-})=\cos t+\sin t\cdot B^{-}\) are \texttt{exp\_of\_sq\_one} and \texttt{exp\_of\_sq\_neg\_one}.

The diagnostic of Section~\ref{sec:motors} is \texttt{Invariant.J5}: the companion scalar \(J=\tfrac12\sum_{a}(\alpha_{a}^{2}-\beta_{a}^{2})\) plus the Minkowski pairing on \(\lambda^{\mu}\). It is unbounded without a constraint on translations (\texttt{J5\_unbounded}).

\subsection{Gravity on a chart}

The diagonal Schwarzschild coframe induces the Schwarzschild metric (\texttt{Gravity.Schwarzschild}, \texttt{inducedMetric\_schwarzschild}). The \((t,r)\) legs are the boost scale factors of the radial rapidity \(\varphi=-\tfrac12\log A\) (\texttt{Gravity.Sandwich}, \texttt{boostScaleFactors\_schwarzschild}), and the finite-angle diagnostic on that slice is \(J=\tfrac12\varphi^{2}\) (\texttt{J\_radialBoostParams}).

\texttt{Gravity.JTDictionary} converts the Weitzenb\"{o}ck scalar on the same chart into
\[
T \;=\; \frac{4}{r^{2}}(\cosh\varphi-1),
\]
which is identical to \(\frac{2}{r^{2}}(1-\sqrt{A})^{2}/\sqrt{A}\) of Section~\ref{sec:gravity} (\texttt{schwarzschild\_T\_eq\_cosh\_rapidity}). Finite angle, field seed, and coframe density remain distinct. The three comparisons of that section are checked in \texttt{Gravity.Identification}: a nonzero finite angle against a flat chart, a pure spatial translation against vanishing \(T\), and the field seed against \(\tfrac12 T\) at \((r_{s},r)=(2,4)\).

\par}

\begin{thebibliography}{99}

\bibitem{Gunn2017}
C.~Gunn,
\textit{Geometric algebras for Euclidean geometry},
Advances in Applied Clifford Algebras \textbf{27} (1), 185--208 (2017).

\bibitem{Lengyel2024}
E.~Lengyel,
\textit{Projective Geometric Algebra Illuminated},
Terathon Software LLC (2024).

\bibitem{DeKeninckDorst2019}
S.~De Keninck and L.~Dorst,
\textit{Projective geometric algebra: A new framework for doing Euclidean geometry},
arXiv:1901.05873 (2019).

\bibitem{GunnDeKeninck2020}
C.~Gunn and S.~De Keninck,
\textit{Geometric algebra and computer graphics},
in: SIGGRAPH 2019 Courses, ACM, New York (2019).
See also: \texttt{https://bivector.net/}.

\bibitem{DorstFontijneMann2007}
L.~Dorst, D.~Fontijne, and S.~Mann,
\textit{Geometric Algebra for Computer Science: An Object-Oriented Approach to Geometry},
Morgan Kaufmann, San Francisco (2007).

\bibitem{Hestenes2015}
D.~Hestenes,
\textit{Space-Time Algebra}, 2nd ed.,
Birkh\"{a}user, Basel (2015).

\bibitem{HestenesSobczyk1984}
D.~Hestenes and G.~Sobczyk,
\textit{Clifford Algebra to Geometric Calculus: A Unified Language for Mathematics and Physics},
Reidel, Dordrecht (1984).

\bibitem{Lounesto2001}
P.~Lounesto,
\textit{Clifford Algebras and Spinors}, 2nd ed.,
London Mathematical Society Lecture Note Series \textbf{286},
Cambridge University Press, Cambridge (2001).

\bibitem{ConwaySmith2003}
J.~H.~Conway and D.~A.~Smith,
\textit{On Quaternions and Octonions: Their Geometry, Arithmetic, and Symmetry},
A K Peters, Natick, MA (2003).

\bibitem{DoranLasenby2003}
C.~Doran and A.~Lasenby,
\textit{Geometric Algebra for Physicists},
Cambridge University Press, Cambridge (2003).

\bibitem{Selig2005}
J.~M.~Selig,
\textit{Geometric Fundamentals of Robotics}, 2nd ed.,
Monographs in Computer Science, Springer, New York (2005).

\bibitem{Weinberg1995}
S.~Weinberg,
\textit{The Quantum Theory of Fields, Volume I: Foundations},
Cambridge University Press, Cambridge (1995).

\bibitem{Tung1985}
W.-K.~Tung,
\textit{Group Theory in Physics},
World Scientific, Singapore (1985).

\bibitem{FultonHarris1991}
W.~Fulton and J.~Harris,
\textit{Representation Theory: A First Course},
Graduate Texts in Mathematics \textbf{129},
Springer, New York (1991).

\bibitem{Humphreys1972}
J.~M.~Humphreys,
\textit{Introduction to Lie Algebras and Representation Theory},
Graduate Texts in Mathematics \textbf{9},
Springer, New York (1972).

\bibitem{maluf2013}
J.~W.~Maluf,
\textit{The teleparallel equivalent of general relativity},
Annalen der Physik \textbf{525} (5--6), 339--357 (2013).

\end{thebibliography}

\end{document}
